% =====================================================================
% 第十一讲 (Exposé XI) 例与补充
% (Exemples et compléments)
% A. Grothendieck
% 原书页码: pp. 219-238
% =====================================================================
\chapter{例与补充}

\section{射影空间，单有理簇}

\begin{proposition}\label{prop:XI.1.1}
设 $k$ 是代数闭域，$X = \mathbf{P}^r_k$ 是 $k$ 上维数 $r$ 的射影空间. 则 $X$ 单连通，i.e. $\pi_1(X) = 0$.
\end{proposition}

对 $r = 0$ 平凡. 若 $r = 1$，需证明：若 $X'$ 是 $X = \mathbf{P}^1_k$ 的非空连通平展覆盖，则 $X' \xrightarrow{\sim} X$. 亏格公式这里给出，若 $g$ 和 $g'$ 是 $X$ 和 $X'$ 的亏格：
\[
1 - g' = d(1 - g),
\]
其中 $d$ 是 $X'$ 在 $X$ 上的次数. 由于 $g = 0$，有 $1 - g' = d$，这要求 $d = 1$，因为 $g' \geqslant 0$，故证明 $X' \xrightarrow{\sim} X$. 当 $r \geqslant 2$ 时对 $r$ 作递归，假设对 $r' < r$ 的 $\mathbf{P}^{r'}$ 单连通. 把它应用于 $\mathbf{P}^r$ 的超平面，用 (\cref{lem:X.2.10})，确实得出 $\mathbf{P}^r$ 单连通. 另一种证明：由 (\cref{cor:X.1.7}) 有 $\pi_1(\mathbf{P}^1 \times \cdots \times \mathbf{P}^1) = \pi_1(\mathbf{P}^1) \times \cdots \times \pi_1(\mathbf{P}^1)$，故 $(\mathbf{P}^1)^r$ 单连通，因为 $\mathbf{P}^1$ 如此，故由基本群的双有理不变性 (\cref{cor:X.3.4})，$\mathbf{P}^r$ 单连通. 这个证明更一般地说明：

\begin{corollary}\label{cor:XI.1.2}
设 $X$ 是代数闭域 $k$ 上紧合的正规概形；若 $X$ 是有理簇，i.e. 整且它的函数域是 $k$ 的纯超越扩张，则 $X$ 单连通.
\end{corollary}

这个结果特别适用于 Grassmann 簇，更一般地适用于 $G/H$，其中 $G$ 是 $k$ 上的连通线性群，$H$ 是含 $G$ 的 Borel 子群 $G$ 的代数子群.

回忆 $k$ 上的\emph{单有理簇}是 $k$ 上紧合且整、函数域 $K$ 包含在 $k$ 的纯超越扩张 $K'$ (在 $K$ 上有限，i.e. 在 $k$ 上有与 $K$ 同样的超越次数) 中的概形，i.e. 存在支配有理映射 $f \colon \mathbf{P}^r_k \to X$，其中 $r = \dim X$. 若 $X$ 正规，由 X 3.4 之前的思考看到，对 $X$ 的每个连通平展覆盖 $X'$ (函数域 $L/K$)，代数 $L \otimes_K K'$ 在 $K'$ 上的模型 $\mathbf{P}^r$ 上非分歧，故由 \ref{prop:XI.1.1} 完全分解，这说明 $L$ 是 $K$-同构于 $K'/K$ 的子扩张. 因此，考虑到 V 8.2，证明：

\begin{corollary}\label{cor:XI.1.3}
代数闭域上正规单有理簇的基本群有限\index[term]{danyoulicu@单有理簇}.
\end{corollary}

(N.B. 注意在"单有理"的定义中，不需要 $K'/K$ 有限).

\begin{remark}\label{rem:XI.1.4}
当然，本节的众所周知. 另一方面，J.-P. Serre 已证明 \cite{exp11-10}：特征零代数闭域上射影光滑单有理簇 $X$ 时，$X$ 单连通. 他的证明是超越的，因为它用 Hodge 对称定理.

\begin{remark}[2003 年 (MR) 增补]
限于特征 $p \geqslant 0$ 的代数闭域 $k$. 在特征 $p > 0$ 时，前面给出的 $k$-簇单有理的定义是弱单有理 $k$-簇的定义，相对的是强单有理 $k$-簇的定义，后者还假设 $K'$ 是 $K$ 的可分扩张.

在维数 2 的情形，存在射影光滑、弱单有理但基本群 (由 \ref{cor:XI.1.3} 有限) 非平凡故非有理的曲面 (T. Shioda，On unirationality of supersingular surfaces，Math. Ann. 225 (1977)，p. 155--159). 相反，强单有理曲面有理：这由 Castelnuovo 有理性判据得出，由 O. Zariski 推广到特征 $p > 0$ (参见 J.-P. Serre，S\'em. Bourbaki no 146，1957 与 \OE uvres--Collected papers，vol. 1，p. 467).

在复数域 $\mathbb{C}$ 上，已知有维数 $\geqslant 3$ 的射影光滑、单有理但非有理簇的例子 (参见 P. Deligne，Vari\'et\'es unirationnelles non rationnelles [d'apr\`es M. Artin et D. Mumford]，S\'em. Bourbaki no 402，1971-72，Lecture Notes vol. 317). 特别在射影空间 $\mathbf{P}^4$ 中的光滑三次超曲面就是这样的情形. 这些例子中有些推广到特征 $> 0$，给出强单有理但非有理的簇 (参见 J.P. Murre，Reduction of the proof of the non-rationality of a non singular cubic threefold to a result of Mumford，Compositio Math. 27 (1973)，p. 63--82).

设 $V$ 是射影正规整 $k$-簇. 称 $V$ 有理连通，若存在整的有限型 $k$-概形 $T$ 和 $k$-态射 $F \colon T \times \mathbf{P}^1 \to V$，使态射：
\begin{equation}
\begin{array}{ccc}
T \times \mathbf{P}^1 \times \mathbf{P}^1 & \longrightarrow & V \times V \\
(t, u, u') & \longmapsto & (F(t, u), F(t, u'))
\end{array}
\tag{1}
\label{eq:XI.1}
\end{equation}
支配. 特别，经过 $V$ 的足够一般的两个有理点有一条有理曲线. 这个术语的合理性在于：若 $V$ 有理连通，可以用有限条有理曲线的链连接两个有理点. 若 $k$ 的特征 $p > 0$，加强前面的定义，要求映射 (1) 一般光滑. 于是得到可分有理连通簇的概念. 这样单有理簇有理连通，且在特征 $p > 0$ 时强单有理簇可分有理连通. J. Koll\'ar 已证明可分有理连通簇有平凡代数基本群 (参见 O. Debarre，Vari\'et\'es rationnellement connexes [d'apr\`es T. Graber，J. Harris，J. Starr et A.J. de Jong]，S\'em. Bourbaki no 906，2001-2002).
\end{remark}
\end{remark}

\section{阿贝尔簇\index[term]{abeiercu@阿贝尔簇}}

设 $k$ 是代数闭域，$A$ 是 $k$ 上的\emph{阿贝尔簇}，i.e. $k$ 上紧合、光滑且连通的群概形，最后 $G$ 是 $k$ 上有限型交换群概形. 用 $\mathrm{Ext}(A, G)$ 记 $A$ 被 $G$ 的交换扩张的类群，用 $H^1(A, G)$ 记 $A$ 上群 $G$ 的主纤维丛的类群 (比较下面第 4 号)，考虑典范同态
\[
\mathrm{Ext}(A, G) \longrightarrow H^1(A, G).
\]
Serre \cite[第七章，定理 5]{exp11-5} 的推理说明它是单射同态，像为 $H^1(A, G)$ 的"本原元素"集，i.e. 满足：
\[
\pi^{*}(\xi) = \mathrm{pr}_1^{*}(\xi) + \mathrm{pr}_2^{*}(\xi)
\]
的元素 $\xi$ 的集，其中 $\mathrm{pr}_i$ 是 $A \times A$ 到 $A$ 的两个投影，$\pi \colon A \times A \to A$ 是 $A$ 的合成律 (N.B. Serre 只在 $G$ 线性且连通、当然在 $k$ 上光滑的情形陈述他的定理，但简化他的推理的第一部分，看到这些限制不必要：只需注意 $A$ 到 $k$ 上有限型群概形 $E$ 的、把单位变成单位的每个态射是群同态，并把它应用于 $A$ 被 $G$ 的扩张 $E$ 在 $A$ 上的截面).

把这个结果应用于 $G$ 是 $k$ 上有限可分群，i.e. 普通有限群，且假设交换的情形. 用 $\pi_1(A \times A) \xrightarrow{\sim} \pi_1(A) \times \pi_1(A)$ (\cref{cor:X.1.7})，并把每个代数概形 $X$ 的 $H^1(X, G)$ 解释为 $\mathrm{Hom}(\pi_1(X), G)$，特别对 $X = A$ 或 $X = A \times A$，看到 $H^1(A, G)$ 的每个类本原，故有同构
\[
\mathrm{Ext}(A, G) \xrightarrow{\sim} H^1(A, G),
\]
换言之，$A$ 的交换结构群 $G$ 的每个在 $A$ 的原点上方带点的主覆盖，唯一地配备把标记点作为原点、且使 $A' \to A$ 成为代数群同态的代数群结构. 特别，若 $A'$ 连通，它也是阿贝尔簇，且与 $A$ 同源.

另一方面，由于函子 $X \mapsto \pi_1(X)$ (从带点代数概形 $X$ 到群) 与乘积交换 (\cref{prop:IX.1.7})，它把第一种范畴中的群变换成群范畴中的群，i.e. 交换群. 故若 $A$ 是阿贝尔簇，$\pi_1(A)$ 是交换群. 故为知道 $\pi_1(A)$，只需知道 $G$ 在有限交换群中变动时的函子 $G \mapsto H^1(A, G) = \mathrm{Hom}(\pi_1(A), G)$. 最后回忆对每个整数 $n > 0$，$A$ 中乘以 $n$ 的同态：
\[
A \xrightarrow{n} A
\]
满射，故核有限，i.e. 是同源，且由此得出每个同源 $A' \to A$ 是前述型同源的商. 由这些及标准推理 (参见例如 \cite{exp11-6}) 得到：

\begin{theorem}\label{thm:XI.2.1}
(Serre-Lang) 设 $A$ 是代数闭域 $k$ 上的阿贝尔簇，对每个整数 $n > 0$ 考虑 $A$ 中乘以 $n$ 的核 ${}_n A$ 的底集普通有限群 $K_n$，最后对每个素数 $\ell$ 令：
\[
T_\ell(A) = \varprojlim_r K_{\ell^r}
\]
且
\[
T_\bullet(A) = \prod_\ell T_\ell(A) = \varprojlim_n K_n
\]
(其中 $m$ 是 $n$ 的倍数 $m = ns$ 时，用乘以 $s$ 把 $K_m$ 送到 $K_n$). 则群 $\pi_1(A)$ 典范同构于 $T_\bullet(A)$，故对每个素数 $\ell$，$\pi_1(A)$ 的 $\ell$-准素分量典范同构于 $T_\ell(A)$.
\end{theorem}

注意这些同构对变动的 $A$ 函子. 模 $T_\ell(A)$ 称为阿贝尔簇 $A$ 的 Tate $\ell$-进模. 它是 $A$ 的加性函子，特别给出 $A$ 的自同态环 $\mathrm{Hom}(A, A)$ 在 $T_\ell(A)$ 上的表示，称为 Weil 的 $\ell$-进表示，它在阿贝尔簇理论中起重要作用 (参见例如 \cite[第七章]{exp11-4}). \cref{thm:XI.2.1} 给出它在 $A$ 的 $\ell$-进同调群 $H_1(A, \mathbb{Z}_\ell) = \pi_1(A)_\ell$ 中的自然表示的解释，后者显然先验地更令人满意，特别从 Lefschetz 公式 \cite[第五章]{exp11-4} 的观点看. 这里回忆 Weil 关于 $T_\ell(A)$ 结构的结果：

a) 若 $n$ 与 $\mathrm{car}(k)$ 互素，则 $K_n$ 是 $\mathbb{Z}/n\mathbb{Z}$ 上秩等于 $2 \dim A$ 的自由模，故若 $\ell$ 是与 $\mathrm{car}(k)$ 不同的素数，$T_\ell(A)$ 是 $\ell$-进整数环 $\mathbb{Z}_\ell$ 上秩等于 $2 \dim A$ 的自由模；

b) 若 $n$ 是 $\mathrm{car}(k) = p$ 的幂，则 $K_n$ 是 $\mathbb{Z}/n\mathbb{Z}$ 上秩 $\nu \leqslant \dim A$ (与 $n$ 无关) 的自由模，故 $T_p(A)$ 是 $p$-进整数环 $\mathbb{Z}_p$ 上秩 $\nu \leqslant \dim A$ 的自由模.

这表明在这里发展的基本群理论中，变动的阿贝尔簇的基本群不随参数正则地变化，它的 $p$-准素分量在对应剩余特征 $p$ 的参数值 $t$ 处可能骤然减小；这个现象最熟知的例子是椭圆曲线. 然而注意，无论 $n$ (是否与特征互素)，$A$ 中乘以 $n$ 的真正核 ${}_n A$ 是 $k$ 上次数 $n^{2g}$ (其中 $g = \dim A$) 的有限群概形，若 $n$ 是 $p = \mathrm{car}(k)$ 的倍数则它在 $k$ 上不可分. 而且，当 $A$ 在阿贝尔簇族中变动，i.e. 在底概形 $S$ 上有阿贝尔概形 $A$ 时，更一般地证明 ${}_n A$ 是 $S$ 上次数 $n^{2g}$ 的有限平坦群概形，即只要考虑核 ${}_n A$ 的无穷小部分，它们对任何 $n$ 都正则地表现. 这提示阿贝尔簇 $A$ 的"真正"基本群是拟代数群 (有限群在 $k$ 上的形式投射极限) $\varprojlim_n {}_n A$，其中代数概形 $X$ 的"真正基本群"应理解为：分类 $X$ 的结构群为 $k$ 上任意有限群 $G$ (不必在 $k$ 上可分) 的主覆盖的拟群. 这样，例如通过 $\mathrm{Hom}(A, A)$ 在 $A$ 的真正基本群的 $p$-准素分量中的表示，比 Serre \cite{exp11-8} 的构造更自然地找回后者用 $\ell \neq p$ 定义的 Weil 特征多项式.

\section{投射锥，Zariski 例子}

为简单仍设 $k$ 是代数闭域，设 $V$ 是 $k$ 上连通的射影 $k$-概形，是 $\mathbf{P}^r_k$ 的闭子概形，若愿意可假设它非奇异. 设 $Y = \widehat{C}_V$ 是 $V$ 的射影\emph{投射锥}\index[term]{toushezhui@投射锥}，$y_0$ 是它的顶点，$X = \widehat{C}_V$ 是与 $\mathcal{O}_V(1)$ 关联的向量丛 $\mathcal{C}_V = \mathbb{V}(\mathcal{O}_V(1))$ 的通常射影闭包，最后
\[
f \colon X \longrightarrow Y
\]
是典范态射，把 $\mathcal{C}_V$ 在 $X$ 上的零截面 $X_0$ 收缩成一点 (EGA II 8.6.4). 由于 $X$ 是 $V$ 上的局部平凡丛，纤维为 $\mathbf{P}^1$ 故单连通，态射 $p \colon X \to V$ 由 (\cref{cor:X.1.4}) 诱导同构：
\[
\pi_1(X) \xrightarrow{\sim} \pi_1(V).
\]
由于 $p$ 诱导同构 $X_0 \to V$，得出 $X$ 的平展覆盖完全分解，当且仅当它在 $X_0$ 上的限制如此. 而 $Y$ 的每个平展覆盖 $Y'$ 的逆像 $X' = X \times_Y Y'$ 是 $X$ 的在纤维 $X_0$ 上完全分解的平展覆盖，故平凡. 由于同态 $\pi_1(X) \to \pi_1(Y)$ 满射 (\cref{cor:IX.3.4})，得出
\[
\pi_1(Y) = (e),
\]
换言之每个射影投射锥单连通. (N.B. 在特征 0 时，取 $Y$ 为仿射投射锥时同样结果成立).

现在假设 $V$ 正则 i.e. 在 $k$ 上光滑；则 $X$ 正则，对 $V$ 的适当射影浸入，找到正规的投射锥 $Y$. 若 $V$ 不单连通，故 $X$ 不单连通，设 $X'$ 是 $X$ 的非平凡连通平展覆盖. 由于 $X$ 在 $y \in Y$ ($y \neq y_0$) 处的纤维约化为一点，看到 $X'$ 在它的诸纤维 (特别在一般纤维) 上的限制平凡；然而 $X$ 不来自 $Y$ 的平展覆盖的逆像，因为 $Y$ 单连通且 $X'$ 将完全分解. 这说明 (X 1.3 和 \ref{rem:XI.1.4}) 在把 $f$ 可分的假设换成更弱的、它的诸纤维是 $k(s)$ 上的可分 (甚至光滑) 代数概形的假设时变假. 同样注意几何纤维 $\overline{X}_y$ ($y \neq y_0$) 的基本群显然约化为 $(e)$，因为这些纤维约化为一点，而 $\pi_1(X_0) \neq e$，故半连续性定理 (\cref{cor:X.2.4}) 对 $f$ 也失效.

最后指出 Zariski 指出的例子，它使这些定理在把 $f$ 可分的假设换成 $f$ 平坦的假设时失效. 设 $f \colon X \to Y$ 是从非奇异射影曲面到有理直线 $Y = \mathbf{P}^1$ 的态射，使 $K = k(X)$ 是 $k(f)$ 的"正则"扩张，i.e. 初等且可分，(i.e. 一般几何纤维连通且可分)，且使除子 $(f) = X_0 - X_\infty$ 是某个除子的 $n$ 倍 (其中 $n$ 是与特征互素的整数). 在任意特征中可以构造这样的例子. 设 $X'$ 是 $X$ 在 $K(f^{1/n})$ 中的正规化，其中 $K = k(X)$ 是 $X$ 的函数域. 由对 $(f)$ 的假设得出 $X'$ 在 $X$ 上平展. 设 $Y'$ 是 $Y$ 在 $k(t)(t^{1/n})$ 中的正规化，它恰在点 $t = 0$ 和 $t = \infty$ 处分歧于 $Y$ 上，且限制 $X'|_{f^{-1}(U)}$ 同构于 $Y'|_U$ 的逆像. 特别，$X'$ 在 $X$ 于 $Y$ 上的一般几何纤维上的限制完全分解. 然而，$X'$ 不同构于 $Y$ 的平展覆盖的逆像，因为立即看到后者必然为 $Y'$，这与 $Y'$ 在 $Y$ 上分歧\footnote{可以注意，从"平展拓扑"(SGA 4 VII) 的观点，在这个例子中 $R^1(f_{\mathrm{ét}})_*(\mathbb{Z}/n\mathbb{Z})$ 在 $S$ 上"不分离".} 矛盾.

下面 (按 Serre) 是实现这个例子条件的简单方式，借鉴 \cite[第 20 号]{exp11-5}：取 $n$ 为 $\geqslant 5$ 的素数，与特征不同，让 $G = \mathbb{Z}/n\mathbb{Z}$ 通过把坐标乘以 $G$ 的四个不同特征 (由于 $n \geqslant 5$ 可能) 作用在 $k^4$ 中\footnote{[2003 年增补：$k^4$ 表示 $k$ 上维数 4 的仿射空间.]}. 则 $G$ 作用在射影空间 $\mathbf{P}^3_k$ 上，$G$ 下唯一的不动点是坐标轴对应的四个点. 曲面 $X'$ 由方程 $x^n + y^n + z^n + t^n = 0$ 定义，在 $k$ 上光滑 (Jacobi 判则)，不含任何不动点，故由于 $G$ 是素阶群，"无不动点"地作用在 $X$ 上，i.e. $X$ 是 $X = X'/G$ 的群 $G$ 的主覆盖. 设 $g = x/y$ 是 $k(X') = K'$ 中的元素，若所选特征是 $\chi_i$，$i = 0, 1, 2, 3$ (其中 $\chi$ 是本原特征)，则它是 $K'$ 在 $K = k(X)$ 上的 Kummer 生成元，设 $f$ 是它的 $n$ 次幂，是 $K$ 的元素. 立即看到 $K'$ 是 $k(g)$ 的正则扩张，这由曲线 (在 $U, T, Z$ 中的齐次方程 $T^n + Z^n + (1 + g^n)U^n = 0$) 在 $k(g)$ 上光滑 (Jacobi 判则) 及已知每条平面曲线连通得出. 另一方面有 $k(f) = K \cap k(g)$，因为第二个成员是 $k(f)$ 的包含在 $k(g)$ (素阶扩张) 中且不同于 $k(g)$ (因为 $g \notin K$) 的扩张. 这蕴含 $K$ 是 $k(f)$ 的正则扩张. 最后 $f$ 在 $X$ 上的除子是某个除子的 $n$ 倍，因为它在 $X'$ 上的逆像是 $g^n$ 的除子，故是 $n$ 倍，且可以下降，因为 $X'$ 在 $X$ 上平展. 若有理映射 $f \colon X \to \mathbf{P}^1$ 是态射，即 $f$ 的零点除子与极点除子不相交，就完成了. 实际上，容易验证 (仍看 $X'$ 上) 这两个除子是 $k$ 上两条在一点 $a$ 横截相交的光滑曲线的 $n$ 倍. 现在把 $X$ 换成在 $a$ 处爆破得到的概形，前面条件 ($\mathrm{div}(f)$ 可被 $n$ 整除，且 $k(X_1) = k(X)$ 是 $k(X)$ 的正则扩张) 仍满足，而且 $f$ 是态射 $X_1 \to \mathbf{P}^1$，故满足所要条件.

\section{上同调正合序列}

设 $S$ 是预概形，故 $X$ 上预概形范畴 $\mathbf{Sch}_{/S}$ 确定，因而其中的群概念也确定，称为 $S$ 上的预概形群，或简称 $S$-群. 为简化叙述并明确思想，此后我们最常限于在 $S$ 上仿射且平坦\footnote{事实上，对下面要做的，用拟仿射假设代替仿射假设就够了，参见第 228 页脚注 (5).} 的群\footnote{参见 SGA 3 IV，特别 §4.}，这足以用于我们考虑的应用. (当然，遇到许多两个假设都不满足的情形). 设 $G$ 是这样的 $S$-群，$P$ 是 $G$ 在其上右作用的 $S$-预概形，这特别蕴含 $\mathbf{Sch}_{/S}$ 中的态射
\[
\pi \colon P \times_S G \longrightarrow P
\]
满足众所周知的公理. 称 $P$ 在 $G$ 下形式主齐性\index[term]{xingshizhuqixing@形式主齐性}，若态射
\[
P \times_S G \longrightarrow P \times_S P
\]
(分量为 $\mathrm{pr}_1$ 和 $\pi$) 是同构；这等价于说对 $\mathbf{Sch}_{/S}$ 的每个对象 $S'$，$P(S') = \mathrm{Hom}_S(S', P)$ 作为带算子群 $G(S') = \mathrm{Hom}_S(S', G)$ 的集，为空或主齐性 (i.e. 为空或同构于 $G(S')$，其上 $G(S')$ 经右平移作用). 称 $P$ 平凡，若 $P$ 同构于 $G$ (其上 $G$ 经右平移作用)，或等价地，若带算子集 $P(S')$ ($G(S')$ 下) 每个都平凡. 例如用著名的归结到集合情形的技巧验证，$P$ 平凡当且仅当它形式主齐性且在 $S$ 上有截面 (后者用范畴术语说 $P$ 在 $\mathbf{Sch}_{/S}$ 的终结对象 $e = S$ 上有截面，i.e. 存在 $e$ 到 $P$ 的态射). 定义比形式\emph{主齐性纤维丛}\index[term]{zhuqixingxianweicong@主齐性纤维丛}更强的 $G$ 下的主齐性纤维丛 $P$ 概念，需先在 $\mathbf{Sch}_{/S}$ 中明确一个用于"下降"的态射集，它起"局部化态射"的作用来"平凡化"纤维丛. 最适当的选择随上下文而变，没有任何一个包含其他所有. 这里方便采用下列定义：

\begin{definition}\label{def:XI.4.1}
设 $G$ 是 $S$-群. 称 $P$ 为 $G$ 下的 (右) 主齐性纤维丛，若 $P$ 是带右 $S$-群 $G$ 的 $S$-预概形，使存在 $S$ 的开集 $U_i$ 的覆盖，以及对每个 $i$ 忠实平坦且拟紧的基变换态射 $S'_i \to U_i$，使 $P' = P \times_S S'$ 是 $G' = G \times_S S'$ 下的平凡带算子预概形，其中 $S'$ 是诸 $S'_i$ 的无交和 $S$-预概形.
\end{definition}

(注意基变换函子 $X \mapsto X' = X \times_S S'$ 左正合，把群变成群，带算子群的对象变成带算子群的对象). 注意 \ref{def:XI.4.1} 随基变换稳定. 还注意：

\begin{proposition}\label{prop:XI.4.2}
设 $G$ 是在 $S$ 上平坦且拟紧的 $S$-群，$P$ 是 $G$ 在其上右作用的 $S$-预概形. 下列条件等价：
\begin{enumerate}[i)]
\item $P$ 是 $G$ 下的主齐性纤维丛；
\item $P$ 在 $G$ 下形式主齐性，且结构态射 $P \to S$ 忠实平坦且拟紧.
\end{enumerate}
\end{proposition}

若 $P$ 在 $G$ 下主齐性，则用 \ref{def:XI.4.1} 的记号，$P'$ 在 $S'$ 上忠实平坦且拟紧 (因为 $G'$ 如此，且 $P'$ 与它 $S'$-同构)，故 $P$ 在 $S$ 上方有同样性质 (对"满射"和"拟紧"，参见 VIII 3.1，对"平坦"是第 VIII 讲的推论中的疏漏). 反之，若 (ii) 满足，取基变换 $S' = P$，它在 $S$ 上确实忠实平坦且拟紧；则 $P'$ 在 $S'$ 上形式主齐性，因为 $P$ 在 $S$ 上如此且基变换函子左正合，另一方面 $P'$ 在 $S'$ 上有截面，即对角截面，故它是平凡主纤维丛，这完成证明.

\begin{corollary}\label{cor:XI.4.3}
若 $G$ 在 $S$ 上仿射且平坦，则 $G$ 下的每个主齐性纤维丛 $P$ 在 $S$ 上仿射且平坦.
\end{corollary}

确实，经基的忠实平坦拟紧扩张后它如此，应用 (\cref{cor:VIII.5.6}).

\cref{def:XI.4.1} 对 $S$ 上平坦仿射 $S$-群有用，因为 (\cref{thm:VIII.2.1})，i.e. \ref{def:XI.4.1} 中考虑的态射 $S' \to S$ 对互相仿射预概形的范畴是有效下降态射. 借助这个事实，下面概述的事实的验证基本上"范畴式"进行\footnote{参见第 226 页脚注 (3) 中的 loc. cit..}. 设 $E$ 是 $S$-群 $G$ 在其上左作用的 $S$-预概形，设 $P$ 是 $G$ 下的 (右) 主齐性纤维丛，想定义"局部"同构于 $E$ 的相伴纤维丛 $E^{(P)}$. 为此，让 $G$ 按律 $(x, y) \mapsto (xg, g^{-1}y)$ 右作用在 $P \times_S E$ 上，该律在集合情形描述此类作用，并经著名技巧推广到范畴. 令 (假定存在)：
\[
E^{(P)} = (P \times_S E)/G
\]
由此看到 $P \times_S E$ 是 $T = E^{(P)}$ 上带右算子群 $G_T = G \times_S T$ 的预概形；为方便，还希望 $P \times_S E$ 是 $T$ 上群 $G_T$ 的主齐性纤维丛. 为验证 $E^{(P)}$ 的存在和前述性质，取\cref{def:XI.4.1} 中的 $S'$，看初始情形在 $S'$ 上的逆像情形. 由于 $P'$ 平凡，i.e. 同构于 $G'_d$，立即看到 $E'^{(P')}$ 存在且有所要的正合性质. 事实上，$E' \times_{S'} P'$ 是 $G'$-同构于 $E' \times_{S'} G'$ 的，故 $E'^{(P')}$ 同构于 $E'$. 而且，带算子空间情形的"相伴纤维丛"的构造与任何基扩张交换，并取种种基扩张 $S'' \rightrightarrows S'$ 和 $S''' \rightrightarrows S''$，其中 $S''$ 和 $S'''$ 是 $S'$ 在 $S$ 上的二重和三重纤维积，看到 $E'^{(P')}$ 配备相对态射 $S' \to S$ 的下降数据，且 $E^{(P)}$ 以所需性质存在当且仅当该下降数据有效；当然 $E^{(P)}$ 正是下降对象. (用 $S' \to S$ 是 $S$-预概形范畴中的下降态射的事实，参见 VIII 5.2). 由此得出 $E$ 在 $S$ 上仿射时相伴纤维丛存在. 把这个构造应用于有 $S$-群同态 $G \to H$，取 $E$ 为 $S$-预概形 $H$ (配备由给定同态得到 $G$ 在 $H$ 上的左作用) 的情形；由于 $H$ 以与 $G$ 在 $H$ 上的作用交换的方式右作用于自身，且 (在 $S$ 上方存在的前提下) 相伴纤维丛的构造与基扩张交换，容易看到 $H$ 将右作用在 $P^{(H)}$ 上，后者按 \ref{def:XI.4.1} 的意义成为 $H$ 下的主齐性纤维丛，且精确地说被与 $P$ 同样的态射 $S' \to S$ 平凡化. 特别，对 $G$ 下的每个主齐性纤维丛 $P$ 和每个 $S$-群同态 $G \to H$ ($H$ 在 $S$ 上仿射)，以对 $(G \to H)$ 函子且与任何基变换 $T \to S$ 相容的方式，关联 $H$ 群的主齐性纤维丛.

\begin{definition}\label{def:XI.4.4}
设 $G$ 是 $S$-预概形. 用 $H^0(S, G)$ 记 $G$ 在 $S$ 上的截面集，当 $G$ 是 $S$-群时把它看作群. 此时用 $H^1(S, G)$ 记 $S$ 上群 $G$ 的主齐性纤维丛的同构类集，把 $H^1(S, G)$ 看作配备对应平凡纤维丛的"标记点"\footnote{这个记号只有在有下降有效性判据时才与一般上同调记号 (SGA 4 V) 相容，后者只在 $G$ 仿射 (或只拟仿射，参见 (\cref{cor:VIII.7.9})) 时几乎不保证.}.
\end{definition}

这样，$H^0(S, G)$ 是 $S$-预概形 $G$ 的、取值在集合范畴的函子. 该函子左正合，更不用说与有限积交换，这确实蕴含它把群变成群，交换群变成交换群. 类似地，借助相伴纤维丛的构造，$H^1(S, G)$ 是仿射 $S$-群 $G$ 的、取值在集合范畴的函子；容易看到该函子与有限积交换. 特别它把第一种范畴中的群，i.e. 交换仿射 $S$-群，变成第二种范畴中的群，甚至交换群 (因为第一种范畴中的群交换). 这样，若 $G$ 是交换仿射 $S$-群，$H^1(S, G)$ 是交换群，交换仿射 $S$-群的同态 $G \to H$ 给出群同态 $H^1(S, G) \to H^1(S, H)$.

为简单，此后限于交换仿射 $S$-群. 设
\[
0 \longrightarrow G' \xrightarrow{u} G \xrightarrow{v} G'' \longrightarrow 0
\]
是这样的群的态射序列，称该序列正合，若 $vu = 0$ (这可以把 $G$ 看作 $G''$ 上带右算子群 $G'_{G''}$ 的预概形) 且 $G$ 是 $G''$ 上群 $G'_{G''} = G' \times_S G''$ 的主齐性纤维丛. 这特别蕴含 $u \colon G' \to G$ 是 $v$ 的核，更不用说蕴含序列
\[
0 \longrightarrow H^0(X, G') \longrightarrow H^0(X, G) \longrightarrow H^0(X, G'')
\]
的正合性. 还蕴含可以定义映射
\[
\partial \colon H^0(X, G'') \longrightarrow H^1(X, G'),
\]
把 $G''$ 在 $S$ 上的每个截面，i.e. 每个 $S$-态射 $f \colon S \to G''$，配给 $S$ 上群 $G' \simeq f^{*}(G'_{G''})$ 的主齐性纤维丛 $P_f$，即 $G''$ 上主齐性纤维丛 $G$ 的逆像. 从 $S$-预概形观点，它正是经 $v \colon G \to G''$ 的、浸入 $f$ 的像 $S$ 的子预概形的逆像，且 $G'$ 在 $P_f$ 上的作用由 $G'$ 在 $G$ 上的右作用诱导. 下列命题的验证也没有除写法以外的困难，留给读者：

\begin{proposition}\label{prop:XI.4.5}
映射 $\partial \colon H^0(X, G'') \to H^1(X, G')$ 是群同态. 下列同态序列正合：
\[
0 \longrightarrow H^0(X, G') \xrightarrow{u_0} H^0(X, G) \xrightarrow{v_0} H^0(X, G'') \xrightarrow{\partial} H^1(X, G') \xrightarrow{u_1} H^1(X, G) \xrightarrow{v_1} H^1(X, G'')
\]
(其中除 $\partial$ 外的同态来自 $H^0$ resp. $H^1$ 的函子律).
\end{proposition}

\begin{remark}\label{rem:XI.4.6}
这里用于研究主齐性纤维丛的观点显然受 Serre \cite{exp11-7} 启发，读者会有兴趣查阅. 当想要一个也适用于在 $S$ 上拟射影的结构 $S$-群 (从而特别包含射影阿贝尔概形) 的形式体系时，有兴趣修改 \ref{def:XI.4.1}，要求 $S'$ 是在覆盖 $S$ 的 $S$ 的开集 $S_i$ 上有限且局部自由的预概形 $S'_i$ 的和. 前面的展开此时成立，特别包括 \ref{prop:XI.4.5}，只要处处把仿射假设换成拟射影假设，并相应解释 $S$-群正合序列的定义. 确实只需把 (\cref{thm:VIII.2.1}) 的引用换成 (\cref{cor:VIII.7.7})：所用态射 $S' \to S$ 对互相拟射影预概形的纤维范畴是有效下降态射. 然而注意这个第二种主齐性纤维丛概念比第一种 \ref{def:XI.4.1} 更严格.
\end{remark}

\refstepcounter{theorem}\noindent\textbf{\thesection.\arabic{theorem}.}\quad 要求 $S$ 被开集 $S_i$ 覆盖，使对每个 $i$，$P|_{S_i}$ 是 $G|_{S_i}$ 下的平凡带算子纤维丛，得到更严格的主齐性纤维丛概念：此时称 $P$ 是\emph{局部平凡主齐性纤维丛}\index[term]{jubupingfanzhuqixingxianweicong@局部平凡主齐性纤维丛}. 这些纤维丛 (对给定的 $G$) 的类构成 $H^1(X, G)$ 的部分，与 $H^1(X, \mathcal{O}_X(G))$ 一一对应，其中 $\mathcal{O}_X(G)$ 是 $G$ 在 $S$ 上的截面的 (通常意义) 层，参见 \cite{exp11-2}. 为使这些 $H^1$ 还给出上同调正合序列 \ref{prop:XI.4.5}，显然必须假设序列 $0 \to G' \to G \to G'' \to 0$ 对这个新上下文在合理意义下正合，i.e. $G$ 是 $G''$ 上群 $G'_{G''}$ 的局部平凡纤维丛；这也意味着 $u \colon G' \to G$ 是 $v \colon G \to G''$ 的核，且 $G$ 在 $G''$ 上局部有截面.

\refstepcounter{theorem}\noindent\textbf{\thesection.\arabic{theorem}.}\quad 引入高上同调群 $H^i(X, G)$ 延续正合序列 \ref{prop:XI.4.5} 显然非常可取. 这可以在"Weil 上同调"框架中做到：考虑 $S$ 上拟紧预概形范畴 $B$，配备忠实平坦拟紧态射集 $M$，称为局部化态射. $S$ 上 (或更好，在 $(B, M)$ 上) 的 Abel "Weil 层"是把和变成积、把
\[
T'' = T' \times_T T' \xrightarrow{\mathrm{pr}_1, \mathrm{pr}_2} T' \xrightarrow{f} T \quad (f \in M)
\]
的序列变成集合的精确图式
\[
F(T) \longrightarrow F(T') \rightrightarrows F(T'')
\]
的、从 $B$ 到 Abel 群范畴的反变函子 $F$. Weil 层构成 Abel 范畴，有正合归纳极限且有一个生成元，故有足够多内射对象 \cite{exp11-1}. 函子 $\Gamma(F) = F(S)$ 的右导出函子记为 $H^i(S, F)$. 另一方面，每个交换 $S$-群显然定义 Weil 层 (\cref{thm:VIII.5.2})，它的 $H^0$ 和 $H^1$ 正是 $H^0(S, G)$ 和 $H^1(S, G)$，从而可以合理方式定义其他 $H^i(S, G)$. 还证明 $S$-群的正合序列定义 Weil 层的正合序列，从而可以找回并延续正合序列 \ref{prop:XI.4.5}\footnote{关于这个观点的系统研究，参见 SGA 4 I 至 IX.}.

\refstepcounter{theorem}\noindent\textbf{\thesection.\arabic{theorem}.}\quad 如 \cite{exp11-2} 那样发展 \ref{prop:XI.4.5} 的非交换变体是适宜的. 关于本节概述的各种上同调概念的适当框架中的系统发展，参见 J. Giraud 的预备工作\footnote{参见 J. Giraud，Cohomologie non ab\'elienne，Springer-Verlag 1971.}.

\section{主纤维丛的特殊情形}

现在假设 $S$ 连通，配备几何点 $a$，从而有分类 $S$ 的平展覆盖的基本群 $\pi_1(S, a)$：$S$ 的平展覆盖范畴等价于 $\pi_1$ 连续作用的有限集范畴. 由此得出 $S$ 上有限平展群概形 $G$ 本质上由普通有限群 $\underline{G}$ (其上 $\pi_1$ 以群自同构连续作用) 确定. $G$ 在其上右作用的 $S$ 的平展覆盖 $P$ 本质上由有限集 $\underline{P}$ (其上 $\pi_1$ 连续地 (左) 作用) 确定，$\underline{G}$ 以与 $\pi_1$ 的作用相容的方式右作用于其上：
\[
s(p \cdot g) = (sp) \cdot (sg) \quad \text{对 } s \in \pi_1, p \in \underline{P}, g \in \underline{G}.
\]
验证 $P$ 按 \ref{def:XI.4.1} 的意义是主齐性纤维丛，当且仅当 $\underline{P}$ 是 $\underline{G}$ 下的主齐性集 (例如用判则 \ref{prop:XI.4.2}). 换言之，$S$ 上群 $G$ 的主齐性纤维丛范畴等价于 $\pi_1$ 连续作用的有限集范畴中群 $\underline{G}$ 的主齐性纤维丛范畴. 特别由此得到对 $G$ 函子的典范双射：
\[
(\ast) \quad H^1(S, G) \xrightarrow{\sim} H^1(\pi_1, \underline{G}),
\]
其中第二个成员表示 $\pi_1$ 作用的 (当然无需说明：连续) 有限集范畴中 $\underline{G}$ 下的主齐性纤维丛的同构类集，该集以熟知方式明确为 1-上链集 $Z^1(\pi_1, \underline{G})$ (满足 $\phi(1) = 1$，$\phi(st) = \phi(s)(s.\phi(t))$) 的商集，除以自然作用其上的群 $\underline{G}$.

一个重要的情形是 $\pi_1$ 平凡作用于 $\underline{G}$，i.e. $G$ 是 $S$ 的完全分解覆盖，同构于 $\underline{G}$ 个 $S$ 的和；此时也写 $H^1(S, \underline{G})$ 代替 $H^1(S, G)$，且该集与 $H^1(\pi_1, \underline{G}) = \mathrm{Hom}(\pi_1, \underline{G})/\underline{G}$ 的内自同构一一对应 ($\ast$). 还注意此时 $S$ 上群 $G$ 的主齐性纤维丛正是 $S$ 的群 $\underline{G}$ 的主覆盖 (\cref{cor:V.2.7})，前述一一对应正是从 $S$ 的群 $\underline{G}$、在 $a$ 上方带点的主覆盖与 $\pi_1(S, a)$ 到 $\underline{G}$ 的连续同态之间的对应 (V 第 5 号末) 推出的对应.

把平展覆盖理论与主纤维丛理论联系起来的兴趣 (已隐含在 A. Weil，G\'en\'eralisation des Fonctions Ab\'eliennes 中，并由 S. Lang 在他的类域论几何理论中明确，参见 Serre \cite{exp11-5})，来自每个在 $S$ 上有限且平展的 $S$-群可以浸入一个在 $S$ 上仿射且光滑、纤维连通、$G$ 交换时也交换的 $S$-群 $H$，使得由正合序列 \ref{prop:XI.4.5} (及其可能的非交换变体)，群 $G$ 的主覆盖的"离散"分类可以用群 $H$ 的主纤维丛的"连续"分类研究，反之亦然. 关于 $G$ 在 $H$ 中浸入的一般构造的想法 (看来在实践中很少用)，参见 \cite[VI 2.8]{exp11-5}. 我们只在下节发展两个重要的、当然经典的特殊情形. 需要一个辅助结果：

\begin{proposition}\label{prop:XI.5.1}
设 $S$ 是预概形，$G$ 是同构于 $\mathrm{GL}(n)_S$ (例如 $G_{m,S}$) 或 $G_{a,S}$ 的 $S$-群，则 $G$ 下的每个主齐性纤维丛局部平凡.
\end{proposition}

明确 $\mathrm{GL}(n)_S$ ($n$ 是 $\geqslant 0$ 的整数) 表示 $S$-预概形 $T$ 上反变函子 $T \mapsto \mathrm{GL}(n, \Gamma(T, \mathcal{O}_T))$ 的 $S$-群，特别 $G_{m,S}$ ("$S$ 上的乘法群") 表示反变函子 $T \mapsto \Gamma(T, \mathcal{O}_T^{*})$，故作为 $S$ 上的预概形同构于 $\mathrm{Spec} \mathcal{O}_S[t, t^{-1}]$，其中 $t$ 是未定元. 类似地 $G_{a,S}$ 表示反变函子 $T \mapsto \Gamma(T, \mathcal{O}_T)$，故作为 $S$-预概形同构于 $\mathrm{Spec}(\mathcal{O}_S[t])$，其中 $t$ 是未定元. 注意经分解论证，\ref{prop:XI.5.1} 重新给出 Rosenlicht 的局部平凡性结果，它关于 $G$ 有"合成序列"、相继因子为这里考虑类型的群的情形. (关于主齐性纤维丛局部平凡性问题的更精细研究，参见 \cite{exp11-7} 和 \cite{exp11-3}).

第一个断言由 $G(T) = \mathrm{Aut}(\mathcal{O}_T^n)$ 及 \ref{def:XI.4.1} 中涉及的态射 $S' \to S$ (i.e. 忠实平坦且拟紧的) 对局部同构于 $\mathcal{O}_T^n$ 的模，i.e. 秩 $n$ 局部自由模，的纤维范畴是\emph{有效下降态射} (\cref{rem:VIII.1.12}) 的事实证明. 第二个类似证明，注意此时 $G(T) = \mathrm{Aut}(E_T)$，其中 $E_T$ 是 $\mathcal{O}_T$ 被 $\mathcal{O}_T$ 的平凡扩张 (其中自同构当然必须保持扩张结构)，且 \ref{def:XI.4.1} 中涉及的态射 $S' \to S$ 对 $\mathcal{O}_T$ 被 $\mathcal{O}_T$ 的扩张的纤维范畴是有效下降态射 (容易由 VIII 1.1 得出)，且这样的扩张自动局部平凡.

\begin{remark}\label{rem:XI.5.2}
注意同样的证明类型适用于辛群 $\mathrm{Sp}(2n)_S$，考虑到：局部同构于 $\mathcal{O}_S^{2n}$ 的模上的交错型，若"非退化" i.e. 定义该模与其对偶的同构，则局部同构于标准型. 正交群的结果相反为假，即使在 $S$ 是域的谱时，因为域上可以有不与标准型同构的二次型. 而且基本上在 \cite{exp11-3} 中证明，群 $\mathrm{GL}$，$\mathrm{Sp}$，$G_a$ 及可分解成此类群的群，是这里考虑的局部平凡性结果几乎唯一的群.
\end{remark}

\begin{corollary}\label{cor:XI.5.3}
有典范双射
\[
H^1(S, \mathrm{GL}(n)_S) \xleftarrow{\sim} H^1(S, \mathrm{GL}(n, \mathcal{O}_S)),
\]
特别
\[
H^1(S, G_{m,S}) \xleftarrow{\sim} H^1(S, \mathcal{O}_S^{*}),
\]
和
\[
H^1(S, G_{a,S}) \xleftarrow{\sim} H^1(S, \mathcal{O}_S),
\]
其中第二个成员表示拓扑空间 $S$ 以通常层为系数的上同调.
\end{corollary}

特别，$H^1(S, \mathrm{GL}(n)_S)$ 等同于 $S$ 上秩 $n$ 局部自由模的同构类集，$H^1(S, G_{a,S})$ 等同于模 $\mathcal{O}_S$ 被它自己的扩张的类集.

\section{应用于主覆盖：Kummer 与 Artin-Schreier 理论}

\begin{proposition}\label{prop:XI.6.1}
设 $S$ 是预概形，$n$ 是 $> 0$ 的整数，$u_n \colon G_{m,S} \to G_{m,S}$ 是 $n$ 次幂同态，${}_n \mu_S$\index[notations]{${}_n\mu_S$} 是它的核. 则 ${}_n \mu_S$ 在 $S$ 上有限且秩 $n$ 局部自由，且当且仅当对每个 $s \in S$，$s$ 的特征与 $n$ 互素时它在 $S$ 上平展. 同态序列
\[
0 \longrightarrow {}_n \mu_S \longrightarrow G_{m,S} \xrightarrow{u_n} G_{m,S} \longrightarrow 0
\]
按第 4 号的意义正合. (称之为 $S$ 上相对整数 $n$ 的 \emph{Kummer 正合序列}\index[term]{Kummer zhenghexulie@Kummer 正合序列}).
\end{proposition}

有
\[
G_m = \mathrm{Spec} \mathcal{O}_S[t, t^{-1}],
\]
$u_n$ 对应仿射 $\mathcal{O}_S$-代数上的同态 $u_n$，由
\[
u_n(t) = t^n
\]
给出，另一方面 $G_{m,S}$ 的单位截面对应 $\mathcal{O}_S$-代数的增广同态，由
\[
\varepsilon(t) = 1
\]
给出，它的核是主理想 $(t - 1)$. 后者经 $u_n$ 的像是主理想 $(1 - t^n)$，得到：
\[
{}_n \mu_S = \mathrm{Spec} \mathcal{O}_S[t]/(1 - t^n),
\]
这特别说明 ${}_n \mu_S$ 在 $S$ 上有限，由 $S$ 上秩 $n$ 自由、以诸 $t^i$ ($0 \leqslant i \leqslant n - 1$) 为基的代数定义. 为使它平展于 $s \in S$，必要且充分条件是约化代数 $k[t]/(1 - t^n)$ (其中 $k = k(s)$，由向 $k$ 形式添加 $n$ 次单位根得到) 在 $k$ 上可分，i.e. $1 - t^n$ 在 $k$ 的代数闭包中的根全部不同，这等价于 $n$ 与特征互素. 最后，证明 \ref{prop:XI.6.1} 中同态序列正合归结到由判则 \ref{prop:XI.4.2} 证明 $u_n$ 忠实平坦. 显然可假设 $S$ 仿射，环为 $A$，故 $G_{m,S}$ 仿射，环为 $B = A[t, t^{-1}]$，只需验证 $u_n$ 使 $B$ 成为 $B$ 上秩 $n$ 的自由模，或等价地 $u_n$ 单射，且 $A[t, t^{-1}]$ 是 $A[t^n, t^{-n}]$ 上秩 $n$ 的自由模. 确实容易验证诸 $t^i$ ($0 \leqslant i \leqslant n - 1$) 构成其一在另一个上的基，这完成证明.

\begin{definition}\label{def:XI.6.2}
称 ${}_n \mu_S$ 为 $S$ 上秩 $n$ 的 Kummer 群，称 $S$ 上秩 $n$ 的 Kummer 主覆盖为 $S$ 上秩 $n$ 的 Kummer 群下的每个主齐性纤维丛.
\end{definition}

这些覆盖的集是群，记为 $H^1(S, {}_n \mu_S)$ 或简记 $H^1(S, {}_n \mu)$. 注意 $S$ 上秩 $n$ 的 Kummer 群的构造与基扩张相容，故 ${}_n \mu_S$ 由绝对 Kummer 群 ${}_n \mu$ 经基扩张自 $\mathrm{Spec}(\mathbb{Z})$ 得到.

用 $(\mathbb{Z}/n\mathbb{Z})_S$ 记由普通有限群 $\mathbb{Z}/n\mathbb{Z}$ 定义的 $S$-群. 若 $G$ 是任意 $S$-群，$(\mathbb{Z}/n\mathbb{Z})_S$ 到 $G$ 的 $S$-群同态 $u$ 与 $G$ 在 $S$ 上的 $n$ 次幂为单位截面的截面一一对应，且与基变换相容，对应 $u$ 的 $(\mathbb{Z}/n\mathbb{Z})_S$ 在 $S$ 上由 $\mathbb{Z}/n\mathbb{Z}$ 的生成元 $1 \mod n\mathbb{Z}$ 定义的截面的像. 这样：

\begin{corollary}\label{cor:XI.6.3}
若 ${}_n \mu_S$ 在 $S$ 上平展，则得到 $S$-群同构 $(\mathbb{Z}/n\mathbb{Z})_S \xrightarrow{\sim} {}_n \mu_S$ 与 $\mathcal{O}_S$ 的在 $S$ 的每个连通分量上精确阶为 $n$ 的截面 (这样的截面称为"$S$ 上的本原 $n$ 次单位根") 之间的一一对应. 故为使 ${}_n \mu_S$ 作为 $S$-群同构于 $(\mathbb{Z}/n\mathbb{Z})_S$，必要且充分条件是它平展于 $S$ 上，i.e. $S$ 的剩余特征与 $n$ 互素，且 $S$ 上存在本原 $n$ 次单位根.
\end{corollary}

这解释了在经典 Kummer 理论中假设基域 (充当 $S$ 的角色) 的特征与 $n$ 互素且含 $n$ 次单位根、并选择本原 $n$ 次单位根所起的作用. 一旦有了概形语言，就不必再纠缠这些假设，应当直接对 $\mu_n$ 而不是 $\mathbb{Z}/n\mathbb{Z}$ 推理. 这样，\ref{prop:XI.6.1}，\ref{prop:XI.4.5} 和 \ref{cor:XI.5.3} 的合取给出 Kummer 主覆盖理论与 Picard 群理论之间的下列一般关系：

\begin{proposition}\label{prop:XI.6.4}
设 $S$ 是预概形，有典范正合序列
\[
0 \to H^0(S, \mu_n) \to H^0(S, \mathcal{O}_S^{*}) \to H^0(S, \mathcal{O}_S^{*}) \to H^1(S, \mu_n) \to H^1(S, \mathcal{O}_S^{*}) \to H^1(S, \mathcal{O}_S^{*}),
\]
从而，令 $H^1(S, \mathcal{O}_S^{*}) = \mathrm{Pic}(S)$，并用 ${}_n A$ 和 $A^n$ 对每个 Abel 群 $A$ 记 $A$ 中乘以 $n$ 的核和余核，有正合序列：
\[
0 \longrightarrow H^0(S, \mathcal{O}_S)^{*n} \longrightarrow H^1(S, \mu_n) \longrightarrow {}_n \mathrm{Pic}(S) \longrightarrow 0.
\]
\end{proposition}

将明确两个重要情形，其中这个正合序列的一个或另一个端项为零：

\begin{corollary}\label{cor:XI.6.5}
假设 ${}_n \mathrm{Pic}(S) = 0$ (例如 $S$ 是局部环的谱，或因子环的谱)，设 $A$ 是环 $H^0(S, \mathcal{O}_S)$. 则有典范同构
\[
H^1(S, \mu_n) \xrightarrow{\sim} A^{*}/A^{*n}.
\]
\end{corollary}

这本质上是 $S$ 是域的谱时的 Kummer 理论经典陈述.

\begin{corollary}\label{cor:XI.6.6}
假设 $H^0(S, \mathcal{O}_S)$ 的每个元素都是 $n$ 次幂，例如 $H^0(S, \mathcal{O}_S)$ 是代数闭域的复合，或 $S$ 在代数闭域 $k$ 上既约且紧合. 则有典范同构
\[
H^1(S, \mu_n) \xrightarrow{\sim} {}_n \mathrm{Pic}(S).
\]
\end{corollary}

特别，当 $S$ 在代数闭域 $k$ 上紧合且连通时，这联系 $S$ 的基本群与 $S$ 在 $k$ 上的 Picard 概形 $P$ 的有限阶点；有同构
\[
\mathrm{Hom}(\pi_1(S), \mathbb{Z}/n\mathbb{Z}) \simeq {}_n P(k)
\]
对与特征互素的 $n$ 成立，它在代数几何中常用. 作为应用，当 $P$ 的连通分量 $P^0$ 是维数 $g$ 的完备群概形时，用第 2 号回顾的结果和 Néron-Severi 挠群的有限性，看到对每个与特征互素的素数 $\ell$，化为阿贝尔群的基本群 $\pi_1(S)$ 的 $\ell$-准素分量是 $\ell$-进整数环 $\mathbb{Z}_\ell$ 上有限型、秩 $2g$ 的模 (而且除至多有限个 $\ell$ 值外自由). 如 Serre 所注意的，这在某些条件下可以用来证明：$X$ 是连通 $S$ 上平坦射影概形时，$X$ 的诸纤维的 Picard 概形都有同样维数，只要应用半连续性定理 (\cref{cor:X.2.3})；Serre 的论证在 $X$ 在 $S$ 上的 Picard 概形存在、且 $X$ 在 $S$ 上的诸纤维的连通 Picard 群是紧合群概形时适用，例如 $X$ 在 $S$ 上的几何纤维正规时 (其中 $X$ 总在 $S$ 上平坦射影)，特别 $X$ 在 $S$ 上光滑射影时.

现在设 $p$ 是素数，假设 $S$ 是特征 $p$ 的预概形，i.e. $p \cdot \mathcal{O}_S = 0$. 则 $\mathcal{O}_S$ 中的 $p$ 次幂同态可加，对应态射 (把 $S$ 换成 $S$ 上变动的 $T$ 得到)：
\[
F \colon G_{a,S} \longrightarrow G_{a,S}
\]
因而是 $S$-群同态，称为 Frobenius 同态 (N.B. 这样的态射对每个由素数域 $\mathbb{Z}/p\mathbb{Z}$ 上的预概形 $G_0$ 经基扩张得到的 $S$-预概形 $G$ 定义，且 $G_0$ 是群概形时该态射是群同态). 令：
\[
\wp = \mathrm{id} - F \colon G_{a,S} \longrightarrow G_{a,S}.
\]
另一方面考虑由普通有限群 $\mathbb{Z}/p\mathbb{Z}$ 定义的 $S$-群 $(\mathbb{Z}/p\mathbb{Z})_S$，已说对每个 $S$-群 $G$，$(\mathbb{Z}/p\mathbb{Z})_S$ 到 $G$ 的 $S$-群同态与 $G$ 在 $S$ 上的 $p$ 次幂为单位截面的截面一一对应. 当 $G = G_{a,S}$ 时，它们对应 $G$ 在 $S$ 上的任意截面. 特别取 $G_{a,S}$ 在 $S$ 上对应环层 $\mathcal{O}_S$ 的单位截面的截面，找到 $S$-群同态
\[
i \colon (\mathbb{Z}/p\mathbb{Z})_S \longrightarrow G_{a,S}
\]

\begin{proposition}\label{prop:XI.6.7}
$S$-群同态序列
\[
0 \longrightarrow (\mathbb{Z}/p\mathbb{Z})_S \longrightarrow G_{a,S} \longrightarrow G_{a,S} \longrightarrow 0
\]
正合 (按第 4 号的意义). (称之为 $S$ 上的 \emph{Artin-Schreier 正合序列}\index[term]{Artin-Schreier zhenghexulie@Artin-Schreier 正合序列}).
\end{proposition}

只需在素数域 $k = \mathbb{Z}/p\mathbb{Z}$ 上证明. 只需注意同态 $\wp^{*} \colon k[t] \to k[t]$ (由 $\wp^{*}(t) = t - t^p$ 定义) 使 $k[t]$ 成为 $k[t]$ 上秩 $p$ 的自由模，确切地说 $k[t]$ 是 $k[s]$ (其中 $s = t - t^p$) 上以诸 $t^i$ ($0 \leqslant i \leqslant p - 1$) 为基的自由模.

由此用 \ref{prop:XI.4.5} 和 \ref{cor:XI.5.3} 得出：

\begin{proposition}\label{prop:XI.6.8}
有典范正合序列：
\[
0 \longrightarrow H^0(S, \mathbb{Z}/p\mathbb{Z}) \longrightarrow H^0(S, \mathcal{O}_S) \longrightarrow H^0(S, \mathcal{O}_S) \longrightarrow H^1(S, \mathbb{Z}/p\mathbb{Z}) \longrightarrow H^1(S, \mathcal{O}_S) \longrightarrow H^1(S, \mathcal{O}_S),
\]
从而有正合序列：
\[
0 \longrightarrow H^0(S, \mathcal{O}_S)/\wp H^0(S, \mathcal{O}_S) \longrightarrow H^1(S, \mathbb{Z}/p\mathbb{Z}) \longrightarrow H^1(S, \mathcal{O}_S)^F \longrightarrow 0,
\]
(其中最后一个项的指数 $F$ 表示端自同态 $F$ 的不变子群，等于 $\wp = \mathrm{id} - F$ 的核).
\end{proposition}

再明确两个极端情形：

\begin{corollary}\label{cor:XI.6.9}
假设 $H^1(S, \mathcal{O}_S)^F = 0$，例如 $S$ 是仿射概形. 则，令 $A = H^0(S, \mathcal{O}_S)$，有典范同构
\[
H^1(S, \mathbb{Z}/p\mathbb{Z}) \xrightarrow{\sim} A/\wp A.
\]
\end{corollary}

这是经典形式的 Artin-Schreier 理论，至少在 $A$ 是域的谱时.

\begin{corollary}\label{cor:XI.6.10}
假设 $\wp H^0(S, \mathcal{O}_S) = H^0(S, \mathcal{O}_S)$，例如 $H^0(S, \mathcal{O}_S)$ 是代数闭域的复合，或 $S$ 在代数闭域上紧合. 则有典范同构：
\[
H^1(S, \mathbb{Z}/p\mathbb{Z}) \xrightarrow{\sim} H^1(S, \mathcal{O}_S)^F
\]
\end{corollary}

\begin{remark}\label{rem:XI.6.11}
最后一个断言归功于 J.-P. Serre \cite{exp11-9}. 也可以对任意 $n$ 的 $n$ 结构群 $\mathbb{Z}/p^n\mathbb{Z}$ 发展类似理论，用 Witt 群概形 $W_n$ 代替 $G_a$，参见 loc. cit.. 注意在特征 $p > 0$ 时，Kummer 理论不再给出关于 $p$ 阶主覆盖的信息，因为 $\mu_p$ 此时是"无穷小"群，i.e. 在基上根式，与 $\mathbb{Z}/p\mathbb{Z}$ 没有直接关系；因此乍看，这些覆盖的理论 (当 $S$ 是代数闭域上的紧合概形以明确思想时) 不再服从 \ref{cor:XI.6.6} 那样的 Picard 概形理论. 然而，若回忆 $\mathrm{Pic}_{S/k}$ 在原点的 Zariski 切空间\footnote{关于 $\mathrm{Pic}_{S/K}$ 的定义，参见 A. Grothendieck，S\'em. Bourbaki N 232 (1962年2月).} 等同于 $H^1(S, \mathcal{O}_S)$，看到知道 $\mathrm{Pic}_{S/k}$ 中乘以 $p$ 的核 ${}_p \mathrm{Pic}_{S/k}$ 的群概形，蕴含知道 $H^1(S, \mathbb{Z}/p\mathbb{Z})$ 以及 $H^1(S, \mu_p)$；注意它也蕴含知道 $H^1(S, \alpha_p)$，其中 $\alpha_p$ 表示素数域上的无穷小群概形，即 $F \colon G_a \to G_a$ 的核 (也可描述为维数 1 的平凡受限 $p$-Lie 代数的受限包络代数的谱)：确实 \ref{prop:XI.4.5} 的正合序列这里给出：
\[
H^1(S, \alpha_p) \simeq \mathrm{Ker}(F \colon H^1(S, \mathcal{O}_S) \longrightarrow H^1(S, \mathcal{O}_S)),
\]
更一般地，用 $\alpha_{p^n}$ 记 $G_a$ 中 $F$ 的 $n$ 次迭代的核，有
\[
H^1(S, \alpha_{p^n}) \simeq \mathrm{Ker}(F^n \colon H^1(S, \mathcal{O}_S) \longrightarrow H^1(S, \mathcal{O}_S)).
\]
事实上，知道 ${}_p \mathrm{Pic}_{S/k}$ 等价于知道所有被 $p$ 零化的有限交换代数群 $G$ 的 $H^1(S, G)$，更一般地，知道 ${}_{p^n} \mathrm{Pic}_{S/k}$ 等价于知道所有被 $p^n$ 零化的有限交换代数群 $G$ 的 $H^1(S, G)$，由下列定理得出，它在所考虑的情形同时涵盖 Kummer 理论和 Artin-Schreier 理论：
设 $G$ 是 $k$ 上有限代数群，$D(G) = \mathrm{Hom}_{k\text{-群}}(G, G_m)$ 是它的 Cartier 对偶 (它的仿射代数由 $G$ 的仿射代数的对偶向量空间承载，i.e. 按 Dieudonn\'e-Cartier 的意义由 $G$ 的超代数承载)，则有典范同构：
\[
(\ast) \quad H^1(S, G) \simeq \mathrm{Hom}_{k\text{-群}}(D(G), \mathrm{Pic}_{S/k}).
\]
(N.B. $S$ 是 $k$ 代数闭上的紧合概形，使 $H^0(S, \mathcal{O}_S) = k$). 这个公式还可以表述为：$S$ 的"真正基本群" (第 2 号提到过) 化为阿贝尔群后同构于诸 $D(P_i)$ 的投射极限，其中 $P_i$ 遍历 $\mathrm{Pic}_{S/k}$ 的有限代数子群，记为 $T^{\bullet}(\mathrm{Pic}_{X/k})$. 当 $S$ 是阿贝尔簇时，在 \ref{thm:XI.2.1} 中已见这个群也同构于"真正"Tate 模 $T_\bullet(S) = \varprojlim_n {}_n S$，且同构 $(\ast)$ 此时写成更醒目的形式
\[
\mathrm{Ext}^1(A, G) \simeq \mathrm{Hom}(D(G), B),
\]
其中 $A$ 是阿贝尔簇，$B$ 是它的对偶，$G$ 是 $k$ 上有限代数群. 刚才指出的结果还可以推广到 $k$ 被任意底预概形代替、以及 $G$ 是有限群以外的其他系数群的情形.
\end{remark}

\section{文献}\printbibliography[heading=none]


